% IEEE conference-format paper. Numbers, tables and the "findings" paragraph are generated by the notebook
% (Section 23) into generated/ - they are \input here, so the text always matches the real run.
\documentclass[conference]{IEEEtran}
\usepackage{cite}
\usepackage{amsmath,amssymb}
\usepackage{graphicx}
\usepackage{booktabs}
\usepackage{adjustbox}
\usepackage{xcolor}
\usepackage{url}
\usepackage{microtype}
\usepackage{flushend}
\renewcommand{\dbltopfraction}{0.9}
\renewcommand{\topfraction}{0.9}
\renewcommand{\bottomfraction}{0.8}
\renewcommand{\textfraction}{0.1}
\renewcommand{\floatpagefraction}{0.7}
\renewcommand{\dblfloatpagefraction}{0.7}
\setcounter{dbltopnumber}{2}
\setlength{\textfloatsep}{8pt plus 2pt minus 2pt}
\setlength{\floatsep}{6pt plus 2pt minus 2pt}
\setlength{\intextsep}{6pt plus 2pt minus 2pt}

% ---------- helpers: the paper still compiles before the notebook has been run ----------
\newcommand{\todo}[1]{\textcolor{red}{[#1]}}
\newcommand{\safeinput}[1]{\IfFileExists{generated/#1}{\input{generated/#1}}{\todo{run notebook Section 23 to create \detokenize{#1}}}}
\newcommand{\safefig}[3]{\IfFileExists{figures/#1.png}{\includegraphics[width=#2]{figures/#1.png}}{\fbox{\parbox{#2}{\centering\todo{missing figure #3}}}}}
\InputIfFileExists{generated/numbers.tex}{}{}
\providecommand{\OffAE}{--}\providecommand{\OffMN}{--}\providecommand{\OffArc}{--}\providecommand{\OffProp}{--}\providecommand{\OffFull}{--}
\providecommand{\TgtAE}{--}\providecommand{\TgtMN}{--}\providecommand{\TgtArc}{--}\providecommand{\TgtProp}{--}\providecommand{\TgtFull}{--}
\providecommand{\SrcAE}{--}\providecommand{\SrcMN}{--}\providecommand{\SrcArc}{--}\providecommand{\SrcProp}{--}\providecommand{\SrcFull}{--}
\providecommand{\OffBestBase}{--}\providecommand{\BestBaseline}{--}\providecommand{\GainOverBest}{--}\providecommand{\GainFullOverBest}{--}
\providecommand{\GainTargetOverBest}{--}\providecommand{\PWilcoxon}{--}\providecommand{\NMachines}{7}\providecommand{\NBoot}{500}\providecommand{\NSeeds}{3}
\providecommand{\AblArc}{--}\providecommand{\AblAda}{--}\providecommand{\AblDann}{--}\providecommand{\AblProp}{--}
\providecommand{\ProbeArc}{--}\providecommand{\ProbeProp}{--}

\begin{document}

\title{DA-ArcNet: Lightweight Domain-Adversarial Angular-Margin Learning for Unsupervised Anomalous Sound Detection of Industrial Machines under Domain Shift}

\author{\IEEEauthorblockN{Ardhanarieswarar}
\IEEEauthorblockA{Roll No.\ 24I461\\
Department of Information Technology\\
PSG College of Technology\\
Coimbatore, India\\
24I461@psgtech.ac.in}
\and
\IEEEauthorblockN{Gokul Kannan}
\IEEEauthorblockA{Roll No.\ 24I463\\
Department of Information Technology\\
PSG College of Technology\\
Coimbatore, India\\
24I463@psgtech.ac.in}
\and
\IEEEauthorblockN{Jafer Sathik B}
\IEEEauthorblockA{Roll No.\ 24I464\\
Department of Information Technology\\
PSG College of Technology\\
Coimbatore, India\\
24I464@psgtech.ac.in}
\and
\IEEEauthorblockN{Madhan M}
\IEEEauthorblockA{Roll No.\ 24I466\\
Department of Information Technology\\
PSG College of Technology\\
Coimbatore, India\\
24I466@psgtech.ac.in}}

\maketitle

\begin{abstract}
Machine sounds are a cheap source of condition information, but faults are rare, so detectors must be trained on normal sounds only, and they must keep working when operating speed, load, microphone position or background noise change (domain shift). We study unsupervised anomalous sound detection under domain shift on MIMII~DG and ToyADMOS2 and propose DA-ArcNet, a small convolutional model that combines (i)~a learnable adaptive clip-mean normalisation layer, (ii)~an angular-margin (ArcFace) classifier over machine section and operating condition, (iii)~a gradient-reversal domain adversary intended to make the embedding uninformative about the source/target domain, and (iv)~k-nearest-neighbour scoring against normal reference embeddings, optionally fused with an autoencoder score. We compare it with three re-implemented baselines (DCASE~2022 autoencoder, DCASE~2022 MobileNetV2, ArcFace-kNN) under the DCASE~2022 protocol (AUC, partial AUC), and add an ablation, cross-dataset transfer in both directions, additive-noise robustness, bootstrap confidence intervals, repeated runs, Wilcoxon tests, Grad-CAM and residual-map explanations, and model-size measurements. Over \NMachines{} machines DA-ArcNet obtains an official score of \OffProp{} (target-domain AUC \TgtProp), and the strongest baseline (\BestBaseline) obtains \OffBestBase; the confidence intervals and tests are reported in Section~\ref{sec:results}.
\end{abstract}

\begin{IEEEkeywords}
anomalous sound detection, domain shift, domain generalisation, machine condition monitoring, ArcFace, domain-adversarial training, explainability
\end{IEEEkeywords}

%--------------------------------------------------------------------
\section{Introduction}
Acoustic monitoring can reveal faults in fans, valves, bearings and other machines before they cause a breakdown. In practice only normal operating sounds are available for training because failures are rare and diverse; the task is therefore \emph{unsupervised} anomalous sound detection (ASD). A second difficulty is \emph{domain shift}: the same machine sounds different when its speed or load changes, when the microphone is moved, or when the factory noise changes, and these changes can alter a recording more than a real fault does~\cite{wilkinghoff2025review}. The DCASE~2022 Challenge Task~2 made this explicit by providing many source-domain but only a few target-domain normal training clips and mixing both domains in the test set without telling the system which is which~\cite{dohi2022dcase}.

The best recent systems rely on large pre-trained feature extractors, adapter fine-tuning or ensembles~\cite{saengthong2024genrep,zheng2024lora,han2025ssam}, which is expensive for deployment on factory edge devices. Others improve the classification loss~\cite{wilkinghoff2024adaproj} or the feature fusion~\cite{liu2022stgram,wilkinghoff2023ssl} but treat domain invariance implicitly. Most studies also evaluate on one challenge dataset at a time.

\textbf{Contributions.} (1)~We design DA-ArcNet, a lightweight network with adaptive clip-mean normalisation and an explicit gradient-reversal domain adversary on top of an angular-margin embedding with kNN scoring (Section~\ref{sec:method}). (2)~We evaluate it against three baselines on both MIMII~DG~\cite{dohi2022mimiidg} and ToyADMOS2~\cite{harada2021toyadmos2} with the official AUC/pAUC protocol, plus an ablation of each component. (3)~We report cross-dataset transfer, noise robustness, statistical evidence (bootstrap intervals, repeated runs, Wilcoxon tests) and explainability (Grad-CAM, residual maps, embedding analysis). (4)~We release the complete code, notebook and this paper's table generator.

%--------------------------------------------------------------------
\section{Related Work}
\textbf{Benchmarks.} The DCASE~2022 task description defines the source/target protocol, the AUC and pAUC metrics and the autoencoder and MobileNetV2 baselines~\cite{dohi2022dcase}; MIMII~DG supplies five machine types with three sections each~\cite{dohi2022mimiidg}; later editions move to a first-shot setting~\cite{dohi2023dcase,nishida2024dcase,nishida2025dcase}. A review of domain generalisation for ASD summarises the challenge systems~\cite{wilkinghoff2025review}.

\textbf{Self-supervised and metric-learning approaches.} Machine-ID or attribute classification is the common pretext task. STgram-MFN fuses spectrogram and temporal features~\cite{liu2022stgram}; FeatEx, a feature-exchange approach, reports state-of-the-art results on DCASE~2023~\cite{wilkinghoff2023ssl}; AdaProj projects embeddings onto class-specific subspaces with an angular margin~\cite{wilkinghoff2024adaproj}; other work uses self-attention over frequency patterns~\cite{zhang2023selfattn}, a predictive autoencoder with contrastive objectives~\cite{zeng2023geco}, spectrum-level contrastive augmentation~\cite{meng2024contrastive}, or addresses the gap between pre-training and ASD fine-tuning~\cite{jiang2024anopatch}.

\textbf{Pre-trained and generative models.} GenRep uses generic features of a large pre-trained extractor with kNN scoring, MemMixup and domain normalisation~\cite{saengthong2024genrep}; LoRA-based fine-tuning of pre-trained audio models is studied in~\cite{zheng2024lora,han2025ssam}, the latter adding a machine-aware group adapter and a dual-level contrastive loss; a diffusion model on log-mel spectrograms with GMM scoring is proposed in~\cite{gao2026udasd}; generative simulation for evaluating ASD systems appears in~\cite{purohit2024mimiigen}.

\textbf{Gap.} Compared with these works, we target a small model in which domain invariance is imposed \emph{explicitly} by an adversary on an angular-margin embedding, and we evaluate across both datasets together with noise robustness and explanations, which the above papers report less often.

%--------------------------------------------------------------------
\section{Problem Formulation}
For a machine section, let $\mathcal{D}=\{x_i\}$ be normal training clips, mostly from the source domain and a few (here 10) from the target domain. We want an anomaly score $s(x)$ such that anomalous clips score higher than normal clips in \emph{both} domains, together with a decision threshold chosen without anomalous data. The domain label is known during training but \emph{not} at test time. Performance is measured per section by $\mathrm{AUC}_{\mathrm{src}}$, $\mathrm{AUC}_{\mathrm{tgt}}$ and the partial AUC $\mathrm{pAUC}$ over $\mathrm{FPR}\in[0,0.1]$ (McClish-standardised) on all test clips; the official score of a machine is the harmonic mean of these values over its sections~\cite{dohi2022dcase}.

%--------------------------------------------------------------------
\section{Proposed Method}\label{sec:method}
\begin{figure*}[!t]
\centering
\safefig{architecture}{0.72\textwidth}{architecture}
\caption{DA-ArcNet. Training (top): normalised log-mel crops pass through AdaCMN and a separable-convolution backbone to an L2-normalised embedding that feeds an ArcFace classifier and, through a gradient-reversal layer, a domain classifier. Inference (bottom): cosine kNN against normal reference embeddings, optional fusion with the autoencoder score, and a gamma-fitted threshold. Orange: components proposed here.}
\label{fig:arch}
\end{figure*}

\subsection{Input and AdaCMN}
Clips are converted to 64-band log-mel spectrograms $X\in\mathbb{R}^{T\times F}$ and standardised with per-band statistics of the training clips. Constant level and microphone or noise offsets are a main source of source/target differences, and clip-mean normalisation removes them but also removes machine-specific level information. \emph{AdaCMN} lets the network choose per band:
\begin{equation}
\tilde X_{t,f}=X_{t,f}-\sigma(a_f)\,\bar X_f,\qquad \bar X_f=\frac1T\sum_{t}X_{t,f},
\end{equation}
with $a_f$ learnable and initialised at $0$ ($\sigma(0)=0.5$).

\subsection{Backbone and angular-margin embedding}
A $3\times3$ convolution, max-pooling and three stages of two depthwise-separable convolutions (32, 64, 128 channels) with batch normalisation and ReLU, followed by global average pooling and a dense layer, give a 128-dimensional embedding $e$ that is L2-normalised. Classes are the operating conditions $y$ of a machine (section and attribute combination); they are available from the file names, not from anomaly labels. The ArcFace loss~\cite{deng2019arcface} with scale $s$ and margin $m$ is
\begin{equation}
\mathcal{L}_{\mathrm{arc}}=-\frac1N\sum_i\log\frac{e^{s\cos(\theta_{i,y_i}+m)}}{e^{s\cos(\theta_{i,y_i}+m)}+\sum_{j\ne y_i}e^{s\cos\theta_{i,j}}} .
\end{equation}

\subsection{Gradient-reversal domain adversary}
Following~\cite{ganin2016dann}, a small domain classifier $g_d$ (source vs.\ target) is attached to $e$ through a gradient-reversal layer (identity forward, gradient multiplied by $-\lambda$ backward):
\begin{equation}
\mathcal{L}=\mathcal{L}_{\mathrm{arc}}+\mathcal{L}_{\mathrm{dom}},
\end{equation}
so that the domain classifier minimises $\mathcal{L}_{\mathrm{dom}}$ while the reversed gradient makes the encoder maximise it,
with $\lambda$ ramped from 0 to $\lambda_{\max}$ during training. Because only 10 target clips exist per section, target clips are oversampled ten times so that the adversary sees both domains.

\subsection{Scoring and decision}
\textbf{kNN score.} Each test clip is split into overlapping 64-frame crops; for each crop the score is $1-\frac1k\sum$ of the $k$ largest cosine similarities to the reference crops of the same section (all normal training clips, $k=3$), averaged over crops. \textbf{Fusion.} The autoencoder reconstruction error $s_{ae}$ (5-frame windows) is combined by standardising $\log s$ with the mean and standard deviation of held-out normal clips of the section and averaging:
$s_{\mathrm{full}}=\exp\big(\tfrac12(z_{\mathrm{prop}}+z_{\mathrm{ae}})\big)$.
\textbf{Threshold.} A gamma distribution is fitted to the scores of held-out normal source clips and its 90th percentile is the threshold, so no anomalous or test data are used.

%--------------------------------------------------------------------
\section{Experimental Setup}
\subsection{Datasets and pre-processing}
\textbf{MIMII DG}~\cite{dohi2022mimiidg}: fan, gearbox, bearing, slider and valve, three sections each. \textbf{ToyADMOS2}~\cite{harada2021toyadmos2}: the full record (169~GB) needs a separate mixing tool, so we use the ToyCar and ToyTrain machines of the DCASE~2022 development set, which are ToyADMOS2 recordings already split into train and test with source and target domains and the same format as MIMII~DG. Per section there are 990 source and 10 target normal training clips, and 50 normal plus 50 anomalous test clips per domain. Audio is 10\,s at 16\,kHz, converted to log-mel spectrograms (FFT 1024, hop 512, 64 bands). The normalisation statistics are fitted on training clips only; 10\,\% of the source normal clips are held out for early stopping, threshold fitting and score scaling and are \emph{not} used as kNN references; training, held-out and test clips are checked to be disjoint in code.

\subsection{Baselines}
(1)~\textbf{AE}: the DCASE~2022 autoencoder (5-frame windows, $4\times128$ hidden units, 8-unit bottleneck), score = reconstruction error. (2)~\textbf{MNv2}: MobileNetV2 ($\alpha=0.5$, trained from scratch) classifying the section, score $=-\log p(\text{section})$~\cite{dohi2022dcase,sandler2018mobilenetv2}. (3)~\textbf{ArcFace-kNN}: the same backbone, loss and kNN scoring as DA-ArcNet without AdaCMN and the adversary, in the style of recent top systems that classify attributes with an angular margin and score by nearest-neighbour distance~\cite{wilkinghoff2024adaproj,saengthong2024genrep}. All baselines are our own re-implementations with the same data, split and training budget; they are not the original authors' systems.

\subsection{Training details}
Implementation in TensorFlow/Keras on a Colab T4 GPU. Crop-based models: 64-frame crops, 4 random crops per clip per epoch, 30 epochs, batch 64, Adam~\cite{kingma2015adam} with cosine decay from $10^{-3}$, $s=30$, $m=0.3$, $\lambda_{\max}=0.3$, SpecAugment-style masking and noise; the checkpoint with the lowest held-out loss is kept. The autoencoder uses Adam, up to 60 epochs with early stopping. The seed of the main run is 42.

\subsection{Evaluation}
Besides AUC/pAUC/official score we report EER, average precision, and at the fixed threshold accuracy, precision, recall, specificity, F1, MCC and confusion matrices. \emph{Statistics:} $\NBoot$ bootstrap resamples (test clips resampled within section, domain and label) give 95\,\% confidence intervals; \NSeeds{} runs with different initialisation and crop sampling are made on two machines; a one-sided paired Wilcoxon signed-rank test over all machine-sections compares DA-ArcNet with each baseline. \emph{Robustness:} white and pink noise at 20, 10, 5 and 0\,dB SNR is added to the test clips only. \emph{Cross-dataset:} embeddings trained on one dataset's normal clips are evaluated on the machines of the other, using only the latter's normal training clips as kNN references. \emph{Ablation:} AdaCMN and the adversary are added separately and jointly to ArcFace-kNN.

%--------------------------------------------------------------------
\section{Results and Analysis}\label{sec:results}
\begin{figure*}[!t]
\centering
\safefig{explainability_maps}{0.74\textwidth}{explainability\_maps}
\caption{Input, autoencoder residual map and Grad-CAM (one median-scored clip per group).}
\label{fig:expl}
\end{figure*}

\subsection{Main comparison}
Table~\ref{tab:main} gives the official score per machine and Table~\ref{tab:domain} the source and target AUC, pAUC and official score over all machines with 95\,\% bootstrap intervals.
\safeinput{tab_main.tex}
\safeinput{tab_domain.tex}
\noindent\safeinput{findings.tex}

DA-ArcNet is best on slider (0.758 against 0.711 for ArcFace-kNN), valve, ToyCar and, on average, over all machines, but it is behind ArcFace-kNN on fan, gearbox and ToyTrain, and on bearing the MobileNetV2 baseline is clearly best (0.663). The mean gain is therefore small (0.638 against 0.633 over the seven machines) and comes from a few machines, which is why the intervals overlap. The MobileNetV2 baseline stayed at chance level (official score 0.500) on four machines (fan, slider, valve and ToyTrain), i.e.\ it did not train successfully there, so its row should be read as a weak baseline and not as evidence about the MobileNetV2 design.

\begin{figure}[t]
\centering
\safefig{metrics_per_machine}{0.58\columnwidth}{metrics\_per\_machine}
\caption{AUC (source, target), pAUC and official score per machine.}
\label{fig:perm}
\end{figure}

\subsection{Threshold-based metrics}
Table~\ref{tab:thr} lists accuracy, precision, recall, specificity, F1 and MCC at the gamma threshold, and Fig.~\ref{fig:cm} the pooled confusion matrices. Because the threshold is fitted on source data only, false alarms on target-domain normal clips are the expected failure mode (error analysis below).
\safeinput{tab_threshold.tex}
\begin{figure}[t]
\centering
\safefig{confusion_matrix_pooled}{0.52\columnwidth}{confusion\_matrix\_pooled}
\caption{Confusion matrices pooled over all machines (rows: models; columns: both domains, source, target).}
\label{fig:cm}
\end{figure}

\subsection{Ablation}
Table~\ref{tab:abl} adds each component to ArcFace-kNN. The official score moves from \AblArc{} (baseline) to \AblAda{} with AdaCMN, \AblDann{} with the adversary and \AblProp{} with both.
\safeinput{tab_ablation.tex}
Only the domain adversary helped on average (0.660 against 0.653), mainly through a higher target AUC (0.704) and pAUC (0.584); the adaptive normalisation layer lowered the source AUC (0.798 to 0.762) and the official score (0.646), and combining both (0.647) did not recover the loss. All differences are 0.013 or less, which is the same size as the run-to-run standard deviation in Table~\ref{tab:seeds} (0.006 to 0.026), so none of the components can be called significant with the present number of runs.

\subsection{Cross-dataset transfer}
\safeinput{tab_crossdata.tex}
Training on MIMII~DG and testing on ToyADMOS2 keeps ArcFace-kNN at 0.631, close to its in-dataset value of 0.633, but lowers DA-ArcNet from 0.638 to 0.590. In the opposite direction both models are near chance (0.521 and 0.542), and DA-ArcNet is higher on gearbox and slider but lower on fan. Transfer between the two datasets is therefore weak, and the domain adversary did not make it better in the MIMII~DG$\rightarrow$ToyADMOS2 direction.

\subsection{Robustness to additive noise}
\safeinput{tab_noise.tex}
Every model that learned something loses accuracy as the SNR falls. From clean to 0\,dB the official score of the autoencoder drops by 0.064, of ArcFace-kNN by 0.115 and of DA-ArcNet by 0.101; DA-ArcNet is the best model at every SNR (0.537 against 0.518 for ArcFace-kNN at 0\,dB and 0.585 against 0.564 at 5\,dB), which is the one consistent advantage we measured. The MobileNetV2 curve is almost flat only because it is close to chance. White and pink noise are averaged in Table~\ref{tab:noise}; the per-type values are in the notebook output (\texttt{noise\_robustness.csv}).

\subsection{Statistical evidence}
Tables~\ref{tab:domain} and~\ref{tab:seeds} give bootstrap intervals and repeated-run statistics. One-sided paired Wilcoxon signed-rank tests over machine-sections (official score) give $p=0.002$ and $0.005$ for DA-ArcNet against the autoencoder and MobileNetV2, but $p=0.354$ against ArcFace-kNN; for DA-ArcNet+AE the values are $0.000$, $0.004$ and $0.840$ (autoencoder, MobileNetV2, ArcFace-kNN). Machine-sections are not independent samples of a larger population, so the $p$-values are descriptive.
\safeinput{tab_seeds.tex}

\subsection{Model size and speed}
DA-ArcNet has 54\,k parameters (ArcFace-kNN 54\,k, autoencoder 187\,k, MobileNetV2 710\,k), i.e.\ about 3.4 times fewer than the autoencoder and 13 times fewer than MobileNetV2. The inference time per 10-s clip measured in the notebook is 6.9\,ms for DA-ArcNet, 6.6\,ms for ArcFace-kNN, 5.3\,ms for the autoencoder and 4.8\,ms for MobileNetV2, and the adversarial branch is used only during training.

\subsection{Explainability and error analysis}
Fig.~\ref{fig:expl} shows the log-mel input, the autoencoder residual map and the Grad-CAM~\cite{selvaraju2017gradcam} of the anomaly score for normal and anomalous source and target clips. A linear probe separates source from target normal clips with balanced accuracy \ProbeArc{} on the ArcFace-kNN embedding and \ProbeProp{} on the DA-ArcNet embedding (0.5 means the domain cannot be recovered). Grad-CAM and residual maps indicate where the models respond; we did not validate them against ground-truth fault locations.
With the threshold from source clips, the fused detector (the one used for the error tables) raises false alarms on almost all target-domain normal clips of ToyCar (100\,\%), bearing (92\,\%), ToyTrain (79\,\%), slider (79\,\%) and fan (78\,\%), against 9--53\,\% on the source domain for the same machines except bearing, where even source clips are rejected (99\,\%). The ten highest-scoring normal clips are all bearing target clips (velocity conditions 24 and 26). The lowest-scoring (most easily missed) anomalies are source-domain anomalies of fan, valve and ToyTrain, and the source miss rates of ToyTrain (76\,\%) and valve (57\,\%) show that some faults leave the learned embedding almost unchanged. This confirms that the threshold, not the ranking quality, is the weak point under domain shift.

%--------------------------------------------------------------------
\section{Discussion and Limitations}
The threshold is fitted on held-out \emph{source} normal clips only, since 10 target clips per section are too few for calibration; threshold-based metrics are therefore weaker on the target domain than AUC. Only ToyCar and ToyTrain represent ToyADMOS2. The baselines are our own re-implementations, so absolute numbers are not comparable with challenge leaderboards. Few seeds were run, and the explanations are qualitative. Future work: per-domain calibration, distilling a pre-trained audio model into the small network, outlier exposure with other machines' sounds, evaluation on the DCASE 2023--2025 first-shot settings and on field data, and INT8 deployment on an edge device.

\section{Conclusion}
We presented DA-ArcNet, a lightweight domain-adversarial angular-margin model for unsupervised anomalous sound detection under domain shift, and evaluated it against three baselines on MIMII~DG and ToyADMOS2 with confidence intervals, repeated runs, an ablation, cross-dataset and noise experiments, explanations and cost measurements. The results support that the model is very small (54\,k parameters) and slightly more robust to additive noise than the ArcFace-kNN baseline, with an official score of 0.638 against 0.633, but they do not show a significant gain over that baseline (Wilcoxon $p=0.35$, overlapping intervals), they do not show that the adversary produced domain-invariant embeddings (probe accuracy 1.00), and the MobileNetV2 baseline failed on several machines. Future work should calibrate thresholds for the target domain, repair the MobileNetV2 baseline and repeat all machines with more seeds.

\section*{Acknowledgment}
The authors thank the Department of Information Technology, PSG College of Technology, and the DCASE organisers.

\clearpage
\bibliographystyle{IEEEtran}
\bibliography{references}
\end{document}
